\subsection{One-pass exact discovery on an endpoint path}
\label{subsec:path-schur-discovery}

Let $P_N=(v_0,v_1,\ldots,v_N)$ be a path with seed $e_{v_0}$.  Write
\begin{equation}
 a_i:=Q_{v_iv_i}=\frac{1+\alpha}{2},
 \qquad
 \beta_i:=-Q_{v_i v_{i-1}}
 =\frac{1-\alpha}{2\sqrt{d_{v_i}d_{v_{i-1}}}}>0
 \quad(1\leq i\leq N),
 \label{eq:path-schur-coefficients}
\end{equation}
and $c_i:=c_{\rho,v_i}$.  Define forward pivots and transformed demands by
\begin{align}
 \delta_0&:=a_0,
 &\eta_0&:=c_0,
 \notag\\
 \delta_i&:=a_i-\frac{\beta_i^2}{\delta_{i-1}},
 &\eta_i&:=c_i+\frac{\beta_i\eta_{i-1}}{\delta_{i-1}}
 \quad(1\leq i\leq N).
 \label{eq:path-forward-schur}
\end{align}
Every $\delta_i$ is positive because it is a pivot of a leading principal
block of the positive definite tridiagonal matrix $Q$.

\begin{theorem}[First nonpositive Schur demand is the exact free boundary]
\label{thm:path-one-pass}
If $\eta_0\leq0$, then $x^\star(\rho)=0$.  Otherwise let
\begin{equation}
 J:=
 \begin{cases}
  \min\{i-1:1\leq i\leq N,\ \eta_i\leq0\},
    &\text{if this set is nonempty},\\
  N,&\text{otherwise}.
 \end{cases}
 \label{eq:path-boundary-index}
\end{equation}
Then
\begin{equation}
 S^\star(\rho)=\{v_0,\ldots,v_J\}.
 \label{eq:path-exact-support}
\end{equation}
The exact solution is recovered from
\begin{align}
 x_{v_J}^\star&=\frac{\eta_J}{\delta_J},
 \notag\\
 x_{v_i}^\star
 &=\frac{\eta_i+\beta_{i+1}x_{v_{i+1}}^\star}{\delta_i}
 \quad(i=J-1,\ldots,0),
 \label{eq:path-back-substitution}
\end{align}
with zero values after $J$.  Discovery, solution, and the terminal KKT check
cost $O(1+\vol(S^\star(\rho)))$ work and storage.
\end{theorem}

\begin{proof}
If $\eta_0=c_0\leq0$, then at $x=0$ the seed residual satisfies
$b_{v_0}\leq\alpha\rho\sqrt{d_{v_0}}$, while every nonseed residual is zero;
the RPPR KKT conditions hold.

Assume $\eta_0>0$.  Forward Gaussian elimination of the shifted Dirichlet
system on a prefix ends at
\begin{equation}
 \delta_i x_{v_i}-\beta_{i+1}x_{v_{i+1}}=\eta_i.
 \label{eq:path-eliminated-row}
\end{equation}
For $i\leq J$, the definition of $J$ gives $\eta_i>0$.  Therefore
\cref{eq:path-back-substitution} is strictly positive and solves
$Q_{SS}x_S=c_{\rho,S}$ on $S=\{v_0,\ldots,v_J\}$.

If $J<N$, the residual at the only outside boundary vertex is
\begin{equation}
 r_{v_{J+1}}=\beta_{J+1}x_{v_J}^\star.
 \label{eq:path-frontier-residual}
\end{equation}
Eliminating the prefix through $J$ shows
\begin{equation}
 \eta_{J+1}
 =c_{J+1}+\beta_{J+1}\frac{\eta_J}{\delta_J}
 =r_{v_{J+1}}-\alpha\rho\sqrt{d_{v_{J+1}}}.
 \label{eq:path-message-is-kkt}
\end{equation}
Thus the stopping test $\eta_{J+1}\leq0$ is exactly the boundary KKT upper
bound.  All more distant residuals are zero.  Apply
\cref{thm:shifted-dirichlet}; the full-support case has no outside boundary.
The forward pass touches every discovered path edge once, and the reverse
pass does the same.  Path degrees are at most two, proving the work bound.
\end{proof}

The theorem is stronger than continuation on this family: it solves the
\emph{final} regularization problem directly, without visiting the
intermediate $\rho_k$ values.  The scalar $\eta_i$ is precisely the causal
message that repeated row pushes fail to preserve.
